\documentclass[a4paper,11pt]{article}
\usepackage[utf8]{inputenc}
\usepackage[T1]{fontenc}
\usepackage[french]{babel}
\usepackage{geometry}
\geometry{left=2cm,right=2cm,top=1.5cm,bottom=2cm}
\usepackage{amsmath,amssymb}
\usepackage{enumitem}
\usepackage{array}
\usepackage{booktabs}
\usepackage{xcolor}
\usepackage{tcolorbox}
\usepackage{fancyhdr}
\usepackage{tikz}
\usepackage{lastpage}

% Couleurs
\definecolor{bleupro}{HTML}{0969da}
\definecolor{orange}{HTML}{d97706}
\definecolor{vertpro}{HTML}{2f855a}
\definecolor{rouge}{HTML}{c53030}
\definecolor{gris}{HTML}{718096}

% En-tête
\pagestyle{fancy}
\fancyhf{}
\lhead{\small CCF Physique-Chimie — Terminale ICCER}
\rhead{\small Fluides \& Courant alternatif}
\cfoot{\thepage/\pageref{LastPage}}
\renewcommand{\headrulewidth}{0.5pt}

% Boîtes
\tcbuselibrary{skins,breakable}

\newtcolorbox{formulairebox}{
  colback=yellow!5,
  colframe=orange,
  title={\textbf{FORMULAIRE — Rappels de cours}},
  fonttitle=\large\bfseries,
  breakable
}

\newtcolorbox{contextebox}{
  colback=blue!3,
  colframe=bleupro,
  leftrule=4pt,
  rightrule=0pt,
  toprule=0pt,
  bottomrule=0pt,
  boxsep=4pt,
  left=8pt
}

\newtcolorbox{reponsebox}[1][4cm]{
  colback=white,
  colframe=gris!40,
  boxrule=0.5pt,
  arc=3pt,
  height=#1,
  valign=top,
  before skip=4pt,
  after skip=8pt
}

\newcommand{\pts}[1]{\hfill\textcolor{gris}{\small\textbf{#1}}}
\newcommand{\comp}[1]{\colorbox{blue!10}{\scriptsize\textbf{#1}}}

\setlength{\parindent}{0pt}
\setlength{\parskip}{6pt}

\begin{document}

% ===== PAGE DE GARDE =====
\begin{center}
{\LARGE\bfseries\color{bleupro} Contrôle en Cours de Formation}\\[8pt]
{\large Physique-Chimie — Terminale Bac Pro ICCER — Groupement 1}\\[12pt]
{\Large\bfseries Pression dans un fluide \& Puissance en régime sinusoïdal}\\[20pt]

\begin{tabular}{|p{6cm}|p{6cm}|}
\hline
\textbf{Durée :} 45 minutes & \textbf{Barème :} /10 points \\
\hline
\textbf{Documents :} Calculatrice autorisée & \textbf{Formulaire :} ci-dessous \\
\hline
\textbf{Nom :} \dotfill & \textbf{Prénom :} \dotfill \\
\hline
\textbf{Classe :} \dotfill & \textbf{Date :} \dotfill \\
\hline
\end{tabular}
\end{center}

\vspace{10pt}

% ===== FORMULAIRE =====
\begin{formulairebox}

\subsection*{Pression dans un fluide}

\begin{tabular}{>{\bfseries}p{4.5cm} p{5.5cm} p{4cm}}
\toprule
\textbf{Grandeur} & \textbf{Formule} & \textbf{Unités} \\
\midrule
Pression & $P = \dfrac{F}{S}$ & Pa, N, m² \\[6pt]
Pression hydrostatique & $\Delta P = \rho \times g \times \Delta h$ & Pa, kg/m³, m/s², m \\[4pt]
Presse hydraulique & $\dfrac{F_1}{S_1} = \dfrac{F_2}{S_2}$ & N, m² \\[6pt]
Conversions & 1 bar $= 10^5$ Pa & \\[4pt]
\bottomrule
\end{tabular}

\smallskip
$\rho_{\text{eau}} = 1\,000$ kg/m³ \quad $g = 9{,}81$ m/s² \quad $P_{\text{atm}} = 1{,}013 \times 10^5$ Pa $= 1{,}013$ bar

\subsection*{Puissance en régime sinusoïdal}

\begin{tabular}{>{\bfseries}p{4.5cm} p{5.5cm} p{4cm}}
\toprule
\textbf{Grandeur} & \textbf{Formule} & \textbf{Unités} \\
\midrule
Puissance active & $P = U \times I \times \cos\varphi$ & W \\[4pt]
Puissance apparente & $S = U \times I$ & VA \\[4pt]
Puissance réactive & $Q = U \times I \times \sin\varphi$ & VAR \\[4pt]
Facteur de puissance & $\cos\varphi = \dfrac{P}{S}$ & (sans unité) \\[6pt]
Énergie & $E = P \times t$ & J (ou kWh) \\[4pt]
\bottomrule
\end{tabular}

\smallskip
\textbf{Triangle des puissances :} $S^2 = P^2 + Q^2$. Un $\cos\varphi$ proche de 1 est souhaitable (moins de pertes réactives).

\end{formulairebox}

\newpage

% ================================================================
\section*{Exercice 1 — Circuit de chauffage d'un immeuble \pts{5 points}}
% ================================================================

\begin{contextebox}
Un installateur thermique intervient sur le circuit de chauffage d'un immeuble de 4 étages. Le circuit est rempli d'eau ($\rho = 1\,000$ kg/m³). La chaufferie est au sous-sol. Le point le plus haut du circuit se trouve au 4\up{e} étage, à une hauteur $h = 12$ m au-dessus de la chaufferie.

La pression au niveau de la chaufferie est $P_{\text{bas}} = 2{,}5$ bar. Le manomètre du vase d'expansion indique la pression relative (pression au-dessus de la pression atmosphérique).
\end{contextebox}

% Schéma
\begin{center}
\begin{tikzpicture}[scale=0.7, every node/.style={font=\small}]
  % Immeuble
  \draw[thick] (0,0) rectangle (4,6);
  \draw[thin] (0,1.5) -- (4,1.5);
  \draw[thin] (0,3) -- (4,3);
  \draw[thin] (0,4.5) -- (4,4.5);
  \node[right] at (4,0.75) {RDC};
  \node[right] at (4,2.25) {1\up{er}};
  \node[right] at (4,3.75) {2\up{e}};
  \node[right] at (4,5.25) {4\up{e}};
  % Chaufferie
  \draw[fill=red!10, thick] (-1,-1) rectangle (5,0);
  \node at (2,-0.5) {\textbf{Chaufferie}};
  % Circuit
  \draw[thick, bleupro] (1,-0.5) -- (1,5.5);
  \draw[thick, bleupro] (3,-0.5) -- (3,5.5);
  \draw[thick, bleupro] (1,5.5) -- (3,5.5);
  % Hauteur
  \draw[<->, thick, rouge] (5.5,0) -- (5.5,6) node[midway, right] {$h = 12$ m};
  % Pressions
  \node[left, bleupro] at (-0.5,-0.5) {$P_{\text{bas}}$};
  \node[left, rouge] at (-0.5,5.5) {$P_{\text{haut}} = \;?$};
\end{tikzpicture}
\end{center}

\textbf{Q1} \comp{APP} — Convertir la pression $P_{\text{bas}} = 2{,}5$ bar en pascals. \pts{1 pt}

\begin{reponsebox}[2.5cm]
\dotfill

\dotfill
\end{reponsebox}

\textbf{Q2} \comp{REA} — Calculer la différence de pression $\Delta P$ entre la chaufferie et le 4\up{e} étage due à la colonne d'eau. \pts{1 pt}

\begin{reponsebox}[3cm]
\dotfill

\dotfill

\dotfill
\end{reponsebox}

\textbf{Q3} \comp{REA} — En déduire la pression $P_{\text{haut}}$ au 4\up{e} étage (en Pa puis en bar). \pts{1 pt}

\begin{reponsebox}[3cm]
\dotfill

\dotfill

\dotfill
\end{reponsebox}

\textbf{Q4} \comp{VAL} — Le technicien sait que la pression minimale au point haut doit être de 0,5 bar pour éviter la cavitation (bulles d'air). La pression obtenue est-elle suffisante ? \pts{1 pt}

\begin{reponsebox}[2.5cm]
\dotfill

\dotfill
\end{reponsebox}

\textbf{Q5} \comp{ANA} — Si l'immeuble avait 8 étages ($h = 24$ m) au lieu de 4, quelle pression minimale faudrait-il en chaufferie pour maintenir 0,5 bar au point haut ? \pts{1 pt}

\begin{reponsebox}[3.5cm]
\dotfill

\dotfill

\dotfill

\dotfill
\end{reponsebox}

\newpage

% ================================================================
\section*{Exercice 2 — Dimensionnement électrique d'une PAC \pts{5 points}}
% ================================================================

\begin{contextebox}
Un technicien installe une pompe à chaleur (PAC) air/eau pour chauffer un immeuble. La PAC est alimentée en monophasé 230 V / 50 Hz. D'après la plaque signalétique :
\begin{itemize}[nosep]
  \item Puissance active : $P = 3\,200$ W
  \item Facteur de puissance : $\cos\varphi = 0{,}85$
\end{itemize}
La PAC fonctionne 10 heures par jour pendant 180 jours (saison de chauffe). Le prix du kWh est 0,20 €.
\end{contextebox}

% Schéma triangle des puissances
\begin{center}
\begin{tikzpicture}[scale=0.6, every node/.style={font=\small}]
  \draw[thick, bleupro] (0,0) -- (7,0) node[midway, below=4pt] {$P$ (W) — active};
  \draw[thick, rouge] (7,0) -- (7,3.5) node[midway, right=4pt] {$Q$ (VAR)};
  \draw[thick, vertpro] (0,0) -- (7,3.5) node[midway, above=4pt, sloped] {$S$ (VA) — apparente};
  % Angle
  \draw (1.5,0) arc[start angle=0, end angle=26.6, radius=1.5];
  \node at (2,0.5) {$\varphi$};
\end{tikzpicture}
\end{center}

\textbf{Q1} \comp{REA} — Calculer la puissance apparente $S$ de la PAC. \pts{1 pt}

\begin{reponsebox}[3cm]
\dotfill

\dotfill

\dotfill
\end{reponsebox}

\textbf{Q2} \comp{REA} — Calculer l'intensité du courant $I$ absorbée par la PAC. \pts{1 pt}

\begin{reponsebox}[3cm]
\dotfill

\dotfill

\dotfill
\end{reponsebox}

\textbf{Q3} \comp{ANA} — Le disjoncteur de la ligne est calibré à 20 A. Est-il correctement dimensionné ? Justifier. \pts{1 pt}

\begin{reponsebox}[2.5cm]
\dotfill

\dotfill
\end{reponsebox}

\textbf{Q4} \comp{REA} — Calculer l'énergie consommée par la PAC sur une saison de chauffe (en kWh). En déduire le coût annuel. \pts{1 pt}

\begin{reponsebox}[3.5cm]
\dotfill

\dotfill

\dotfill

\dotfill
\end{reponsebox}

\textbf{Q5} \comp{COM} — Le propriétaire demande pourquoi sa facture est élevée. Le technicien remarque que le $\cos\varphi$ est de 0,85 au lieu de 0,95 (valeur idéale). Expliquer en quoi un mauvais facteur de puissance augmente l'intensité appelée, et pourquoi c'est un problème. \pts{1 pt}

\begin{reponsebox}[4cm]
\dotfill

\dotfill

\dotfill

\dotfill

\dotfill
\end{reponsebox}

\vfill

\begin{center}
\begin{tabular}{|c|c|c|c|}
\hline
\textbf{Exercice} & \textbf{Thème} & \textbf{Questions} & \textbf{Points} \\
\hline
1 & Pression dans un fluide (circuit chauffage) & Q1 à Q5 & /5 \\
\hline
2 & Puissance en régime sinusoïdal (PAC) & Q1 à Q5 & /5 \\
\hline
\multicolumn{3}{|r|}{\textbf{Total}} & \textbf{/10} \\
\hline
\end{tabular}
\end{center}

\end{document}
